\slajdid{07-s014}
\begin{frame}[label=07-s014]
\gusttekst\setlength{\parskip}{4pt}
\textbf{2. Jezik za unos dizajna.} Baš kao što programski jezici visokog nivoa dozvoljavaju da se složeni koncepti dizajna izraze kao kompjuterski program VHDL omogućava da se ponašanje složenih elektronskih kola prikaže u sistemu dizajna za automatsku sintezu kola ili za simulaciju sistema. Kao Pascal, C i C++, VHDL uključuje funkcije korisne za tehnike strukturiranog dizajna i nudi bogat skup funkcija kontrole i predstavljanja podataka. Za razliku od ovih drugih programskih jezika, VHDL pruža funkcije koje omogućavaju opisivanje \textcolor{DEcrvena}{istovremenih događaja}. Ovo je važno jer je hardver koji se opisuje korišćenjem VHDL-a inherentno istovremen u svom radu. Korisnicima PLD programskih jezika kao što su PALASM, ABEL, CUPL i drugi će istovremene karakteristike VHDL-a biti prilično poznate. Oni koji su programirali samo koristeći softverske programske jezike će, međutim, imati neke nove koncepte za razumevanje. Konkuretno programiranje. Konkuretno izvršavanje programa. \par\vspace{4mm}
\textbf{3. Jezik za testiranje.} Jedan od najvažnijih (i nedovoljno iskorišćenih) aspekata VHDL-a je njegova sposobnost da prihvati specifikacije performansi za kolo, u obliku koji se obično naziva testbench. Testbench su VHDL opisi stimulusa kola i odgovarajućih očekivanih izlaza koji potvrđuju ponašanje kola tokom vremena. Testbench treba da budu sastavni deo svakog VHDL projekta i treba da se kreiraju paralelno sa drugim opisima kola.
\end{frame}
